% !TeX spellcheck = en_US
% !TeX encoding = UTF-8
\chapter{Task 4: DDoS Filtering with eBPF/XDP}\label{ch:task4}

Detection alone does not stop an attack: once a flood is recognised, its packets must be discarded before they consume the victim's resources, and more cheaply than they are produced. We first summarise eBPF, XDP, their use cases and tools, and then present our per-source-IP rate limiter, the test bed and the measured results.

\section{eBPF and XDP}\label{sec:t4-ebpf}\label{sec:t4-xdp}
The \emph{extended Berkeley Packet Filter} (eBPF) lets small, sandboxed programs run inside the Linux kernel without changing the kernel source or loading a module \autocite{kdoc_bpf,ebpf_what}. Programs are written in restricted C and compiled by clang (\texttt{-target bpf}) to a 64-bit bytecode (11 registers, 512-byte stack) that the kernel JIT-compiles to native code. Before loading, the \textbf{verifier} simulates every path: it checks termination (only bounded loops), that every memory access is in bounds (each packet pointer must be compared with \texttt{data\_end} before use), that no kernel pointer leaks and that only allowed helpers are called \autocite{kdoc_verifier}. State, configuration and statistics live in \textbf{maps}, key/value structures shared with user space (hash, array, LRU hash, per-CPU variants, ring buffer, LPM trie, \dots) \autocite{kdoc_maps,kdoc_maphash}. The program type decides the hook: XDP and tc for networking, kprobes and tracepoints for tracing, cgroup and LSM hooks for security. With the BPF Type Format \autocite{kdoc_btf}, libbpf relocates field offsets at load time (\enquote{Compile Once -- Run Everywhere}) \autocite{nakryiko_core}; our program uses a \texttt{vmlinux.h} generated from the kernel's BTF and a libbpf skeleton.

\textbf{XDP} (eXpress Data Path) is an eBPF hook at the earliest point of the receive path: in the NIC driver, right after DMA and \emph{before} the kernel allocates a socket buffer (\texttt{sk\_buff}) \autocite{hoiland2018xdp,ebpfdocs_xdp}. The program receives a minimal context (\texttt{struct xdp\_md}: frame start/end, ingress interface, RX queue) and returns a verdict: \texttt{XDP\_DROP} discards the frame, \texttt{XDP\_PASS} hands it to the normal stack, \texttt{XDP\_TX} sends it back out of the same NIC, \texttt{XDP\_REDIRECT} to another NIC, CPU or AF\_XDP socket, and \texttt{XDP\_ABORTED} signals an error. Because no \texttt{sk\_buff} exists yet, a drop is very cheap: Høiland-Jørgensen et al.\ report tens of millions of dropped packets per second on one core and performance close to kernel bypass (DPDK), while the kernel keeps control of the NIC \autocite{hoiland2018xdp}; in Cloudflare's measurements, the same drop rule was fastest at the XDP layer, ahead of the iptables \texttt{raw} table and tc \autocite{cf_drop10m}. iptables/nftables remain the tool for stateful, policy-rich firewalling and tc-BPF for egress; XDP is the place for cheap, early decisions on huge packet volumes, i.e.\ the DDoS case. Programs attach in three modes \autocite{ebpfdocs_xdp,cilium_bpfref}: \textbf{native (driver)} mode in the driver's NAPI poll loop (the fast path; most 10--100\,Gbit/s drivers, \texttt{veth}, \texttt{virtio\_net}), \textbf{generic (skb)} mode, a driver-independent fallback after \texttt{sk\_buff} allocation that suits testing but loses most of the speed, and \textbf{offloaded} mode on a SmartNIC.

\section{Use Cases and Tools}\label{sec:t4-usecases}\label{sec:t4-tools}
\textbf{DDoS mitigation} is the classic use case: Cloudflare's \emph{Gatebot} detects attacks and generates rules that \emph{L4Drop} compiles into XDP programs on every edge server \autocite{cf_gatebot,cf_l4drop}, and \emph{flowtrackd} adds stateful TCP tracking \autocite{cf_flowtrackd}. In \textbf{load balancing}, Facebook's \emph{Katran} uses \texttt{XDP\_TX} with IP-in-IP encapsulation \autocite{meta_katran,katran_github}, as does Cloudflare's \emph{Unimog} \autocite{cf_unimog}. For \textbf{firewalling and policy}, \emph{Cilium} uses XDP for a CIDR pre-filter and its kube-proxy replacement in Kubernetes \autocite{cilium_kpr,cilium_bpfref}. For \textbf{monitoring}, XDP programs count, sample and capture packets at line rate (e.g.\ \texttt{xdpdump}). As \textbf{IDS/IPS acceleration}, Suricata bypasses flows it has finished inspecting \autocite{suricata_ebpf}. \textbf{AF\_XDP} redirects selected packets zero-copy to a user-space socket \autocite{kdoc_afxdp}. \Cref{tab:t4-tools} surveys the tools we looked at: building blocks for writing and loading programs, complete systems that use XDP internally, and eBPF-based security tools. For our own program we used libbpf, bpftool and iproute2; Task~5 covers xdp-tools.

\begin{table}[htbp]
	\centering\footnotesize
	\caption{Survey of eBPF/XDP tools.}\label{tab:t4-tools}
	\begin{tabularx}{\linewidth}{>{\raggedright\arraybackslash}p{4.3cm}>{\raggedright\arraybackslash}p{2.2cm}X}
		\toprule
		Tool & Type & Description \\
		\midrule
		xdp-tools \autocite{xdptools_github} & XDP utilities & libxdp (multi-program dispatcher), \texttt{xdp-loader}, \texttt{xdp-filter}, \texttt{xdpdump}, \texttt{xdp-bench} (\cref{ch:task5}). \\
		xdp-tutorial \autocite{xdp_tutorial} & Teaching & Lessons from basic programs to AF\_XDP; basis of our program's structure. \\
		libbpf \autocite{libbpf_github}, bpftool \autocite{kdoc_bpf} & Library, CLI & Reference loader (CO-RE, skeletons); inspect/update maps, generate \texttt{vmlinux.h}. \\
		iproute2 (\texttt{ip link \dots\ xdp obj}) & CLI & Attaches an XDP object without a custom loader. \\
		BCC \autocite{bcc_github}, bpftrace \autocite{bpftrace_github} & Tracing & Python front-end with 100+ tools; awk-like tracing one-liners. \\
		Cilium \autocite{cilium_kpr}, Katran \autocite{katran_github} & Networking & Kubernetes networking with XDP acceleration; XDP L4 load balancer. \\
		Gatebot, L4Drop, flowtrackd \autocite{cf_gatebot,cf_l4drop,cf_flowtrackd} & DDoS pipeline & Detection and rule generation, XDP enforcement, stateful TCP. \\
		Suricata XDP bypass \autocite{suricata_ebpf} & IDS/IPS & Early drop/bypass of inspected flows. \\
		Falco \autocite{falco}, Tetragon \autocite{tetragon} & Runtime security & System-call/process observability and enforcement (not packets). \\
		Polycube \autocite{polycube_github} & NFV framework & Chains eBPF/XDP network functions (router, firewall, DDoS mitigator). \\
		\bottomrule
	\end{tabularx}
\end{table}

\section{Our Filter: a Per-Source-IP Rate Limiter}\label{sec:t4-design}
A volumetric flood such as a SYN flood sends far more packets per second from a source than any legitimate client. Our XDP program therefore counts packets per IPv4 source address in a fixed time window and drops everything above a configurable limit. Limit, window, allow list and block list live in BPF maps and can be changed at runtime without recompiling or reloading. \Cref{fig:t4-arch} shows the architecture.

\begin{figure}[htbp]
	\centering
	\includegraphics[width=0.72\linewidth]{xdp_pipeline.png}
	\caption{Architecture of the XDP rate limiter: per-packet decision path in the kernel (centre), BPF maps shared with user space (right), and the user-space components (top), including the planned interface for an ML classifier.}\label{fig:t4-arch}
\end{figure}

\paragraph{Decision path.} For every frame, the program (\cref{lst:t4-prog}) parses Ethernet and an optional 802.1Q/802.1ad VLAN tag, then the IPv4 header. Every header access is preceded by a bounds check against \texttt{data\_end}, without which the verifier rejects the program. Non-IPv4 frames (ARP, IPv6) are passed and counted. For IPv4, the source address is first looked up in the block list (drop) and the allow list (pass). Then the fixed-window counter runs: the \texttt{rl\_state} entry of the source stores the window start (from \texttt{bpf\_ktime\_get\_ns()}) and the packet count. When the window has expired, it is reset; then the count is incremented atomically, and a packet whose count exceeds the limit gets \texttt{XDP\_DROP}. Limit and window length are read from \texttt{rl\_config} on every packet, so a change by the operator takes effect on the next packet. A limit of 0 disables the limiter.

\lstinputlisting[style=bpfc,float=tbp,firstline=125,lastline=160,caption={Block/allow lists and fixed-window counter in \texttt{xdp\_ratelimit.bpf.c} (Ethernet/VLAN/IPv4 parsing above it omitted).},label=lst:t4-prog]{../xdp/ratelimit/xdp_ratelimit.bpf.c}

\paragraph{Map choices.} \Cref{tab:t4-maps} lists the five maps. The key decision concerns \texttt{rl\_state}. It is an \texttt{LRU\_HASH}, which can never become \enquote{full}: at capacity the least recently used source is evicted, whereas a plain hash map would refuse new entries, which an attacker could exploit. It is \emph{shared}, not per-CPU, on purpose: with receive-side scaling, one source's packets can be spread over several CPUs, and per-CPU counters would raise the effective limit to $\text{limit} \times \#\text{CPUs}$. The price is an atomic \texttt{fetch\_add} on a shared cache line (which requires the BPF v3 instruction set, \texttt{-mcpu=v3}) and a benign race on the window reset, in which at worst a few extra packets pass. The global statistics, which every packet writes, use a per-CPU array instead, so that the hot path needs no atomics there; user space sums the per-CPU slots.

\begin{table}[htbp]
	\centering\small
	\caption{BPF maps of the rate limiter.}\label{tab:t4-maps}
	\begin{tabularx}{\linewidth}{llrX}
		\toprule
		Map & Type & Entries & Key $\to$ value and purpose \\
		\midrule
		\texttt{rl\_config} & \texttt{ARRAY} & 1 & $0 \to \{$limit, window\_ns$\}$; runtime configuration, cheapest possible lookup \\
		\texttt{rl\_state} & \texttt{LRU\_HASH} & 65,536 & source IPv4 $\to \{$win\_start, count, total, dropped$\}$; per-source window state \\
		\texttt{rl\_stats} & \texttt{PERCPU\_ARRAY} & 4 & pass / drop-rate / drop-block / non-IPv4 counters \\
		\texttt{rl\_allow} & \texttt{HASH} & 1,024 & source IPv4 $\to 1$; never rate-limited \\
		\texttt{rl\_block} & \texttt{HASH} & 65,536 & source IPv4 $\to 1$; always dropped (interface for an IDS/ML model) \\
		\bottomrule
	\end{tabularx}
\end{table}

\paragraph{Loader and build.} The user-space loader \texttt{xdp\_ratelimit\_user} (about 210 lines of C) uses the libbpf skeleton. It loads the object (the verifier runs inside \texttt{open\_and\_load}), writes the configuration and the allow/block lists \emph{before} attaching, so that the first packet is already governed, attaches in generic or native mode with \texttt{bpf\_xdp\_attach()}, prints per-second counters and the top sources by drops, and detaches on SIGINT/SIGTERM or after a given duration. Alternatively, the object can be attached with iproute2 alone and configured with \texttt{bpftool map update}. The \texttt{Makefile} generates \texttt{vmlinux.h} from the kernel's BTF (\texttt{bpftool btf dump file /sys/kernel/btf/vmlinux format c}), compiles the BPF object (\texttt{clang -O2 -g -target bpf -mcpu=v3}), generates the skeleton (\texttt{bpftool gen skeleton}) and links the loader with \texttt{-lbpf -lelf -lz}. We built it with clang 21.1.8 and bpftool v7.7.0 on kernel 7.0.0-43-generic. The verifier accepted the program: \texttt{bpftool prog show} reports it as loaded and JIT-compiled (2,208 bytes of translated bytecode, 1,257 bytes of native code).



\section{Test Bed and Demo}\label{sec:t4-demo}
To test the filter without touching the host's interfaces, the script \texttt{run\_demo.sh} builds four network namespaces (\cref{fig:t4-topo}): an attacker (10.10.0.2), a benign client (10.10.0.3) and a victim (10.10.0.1), connected through veth pairs to a Linux bridge in a fourth namespace. The rate limiter is attached to the victim's interface \texttt{vic0}. The attacker runs a SYN flood to port~80 (\texttt{hping3 -S -p 80 --flood}) for 5\,s, while the benign client pings the victim five times per second. We count the TCP segments that reach the victim's IP stack, i.e.\ that passed XDP, with the kernel's own \texttt{TcpInSegs} counter of the victim namespace, read before and after each round. The default limit is 100 packets per second and source.

\begin{figure}[htbp]
	\centering
	\begin{tikzpicture}[
		ns/.style={draw, rounded corners, fill=gray!8, minimum width=3.6cm, minimum height=1.25cm, align=center, font=\small},
		br/.style={draw, rounded corners, fill=blue!6, minimum width=2.6cm, minimum height=2.9cm, align=center, font=\small},
		lnk/.style={thick},
		node distance=1.2cm]
		\node[br] (lan) {\texttt{rl-lan}\\[2pt] bridge \texttt{br0}};
		\node[ns, left=1.6cm of lan, yshift=0.85cm] (att) {\texttt{rl-attacker}\\ 10.10.0.2 (\texttt{att0})\\ \texttt{hping3 -S --flood}};
		\node[ns, left=1.6cm of lan, yshift=-0.85cm] (ben) {\texttt{rl-benign}\\ 10.10.0.3 (\texttt{ben0})\\ \texttt{ping -i 0.2}};
		\node[ns, right=1.9cm of lan, minimum width=4.2cm] (vic) {\texttt{rl-victim}\\ 10.10.0.1 (\texttt{vic0})\\ TCP stack, \texttt{TcpInSegs}};
		\draw[lnk] (att.east) -- node[above, font=\scriptsize]{veth} (lan.west |- att.east);
		\draw[lnk] (ben.east) -- node[below, font=\scriptsize]{veth} (lan.west |- ben.east);
		\draw[lnk] (lan.east) -- node[above, font=\scriptsize]{veth} (vic.west);
		\node[draw=orange!80!black, fill=orange!15, thick, font=\scriptsize\bfseries, align=center, inner sep=2pt] at ([xshift=-1pt]vic.west) [anchor=east] {XDP};
		\node[below=0.15cm of vic, font=\scriptsize, align=center] {\texttt{xdp\_ratelimit} on \texttt{vic0}\\ (generic or native mode)};
	\end{tikzpicture}
	\caption{Namespace test bed of \texttt{run\_demo.sh}. The XDP program filters all traffic that enters the victim namespace.}\label{fig:t4-topo}
\end{figure}

The demo runs labelled phases, all logged with timestamps to \texttt{logs/demo.log}: (0)~topology and connectivity check; (1)~baseline flood \emph{without} XDP; (2)~attach in generic (skb) mode with the libbpf loader; (3)~benign ping alone; (4)~flood and benign ping in parallel with the limiter active (the main result); (5)~map inspection with \texttt{bpftool map dump}; (6)~runtime control: raise the limit to 1,000 via \texttt{bpftool map update}, restore it, and block-list the attacker; (7)~the same program in native (driver) mode, which veth supports; (8)~loading with iproute2 only (\texttt{ip link set dev vic0 xdpgeneric obj \dots}), configuration with bpftool and a 3\,s flood.

\section{Results}\label{sec:t4-results}
\Cref{tab:t4-results} summarises the attack rounds. Without XDP, all 331,744 flood packets reached the victim's TCP stack (the kernel counted 331,745 TCP segments; one-packet differences between the hping3 summary and the kernel counters are within the counting granularity). With the rate limiter in generic mode, the attacker sent 679,608 packets in 5\,s, and exactly 500 TCP segments reached the victim, i.e.\ 100 per second for 5 seconds, as configured. Native mode gave the same result (613,271 packets sent, 500 received; \texttt{rl\_state} counted 613,272 for the attacker, the same counting granularity), and so did the iproute2-loaded object over 3\,s (383,335 sent, 300 received). In all rounds, the benign client lost no ping. The attacker could send about twice as many packets with the filter active as in the baseline; we attribute this to the CPU time that the victim's stack spent on processing (and answering) every SYN in the baseline, since all namespaces share the same CPUs.

\begin{table}[htbp]
	\centering\small
	\caption{Results of the attack rounds (limit 100 packets/s per source; values from \texttt{logs/demo.log}, \texttt{hping3\_*.log}, \texttt{victim\_*.log} and \texttt{benign\_*.log}).}\label{tab:t4-results}
	\begin{tabular}{llrrrc}
		\toprule
		Round & XDP mode / loader & Flood & Packets sent & Segments at victim & Benign ping loss \\
		\midrule
		Baseline & none & 5\,s & 331,744 & 331,745 & 0\,\% (25/25) \\
		Phase 4 & generic, libbpf loader & 5\,s & 679,608 & 500 & 0\,\% (25/25) \\
		Phase 7 & native, libbpf loader & 5\,s & 613,271 & 500 & 0\,\% (25/25) \\
		Phase 8 & generic, iproute2 + bpftool & 3\,s & 383,335 & 300 & 0\,\% (15/15) \\
		\bottomrule
	\end{tabular}
\end{table}

\paragraph{Drop rate and maps.} The loader's per-second statistics show between 112,135 and 140,096 rate-limit drops per second in generic mode and between 123,321 and 130,299 in native mode (full seconds of the flood), with a steady 104--106 passed packets per second (100 from the attacker plus the benign pings). These rates are limited by what hping3 could offer on the same machine, not by the filter, so they are not a throughput measurement of XDP; with a single sender, generic and native mode cannot be distinguished. The map dump after phase~4 (\cref{fig:ss-demo_phase5_maps.png}) shows the expected state: \texttt{rl\_state} holds the attacker 10.10.0.2 with \texttt{total} 679,608 and \texttt{dropped} 679,108 (exactly 500 passed) and the benign client 10.10.0.3 with 40 packets and none dropped; the per-CPU \texttt{rl\_stats} slots show that the drops were spread over several CPUs (three in phase~4), which justifies the shared state map.

\paragraph{Runtime control and block list.} In phase~6 we raised the limit to 1,000 packets/s by writing the little-endian value into \texttt{rl\_config} with \texttt{bpftool map update}; the following dump confirmed the new value without any reload (the flood had already ended, so this phase demonstrates the control path, not its effect under load). We then restored the limit and inserted the attacker into \texttt{rl\_block}. Afterwards, the attacker's pings had 100\,\% loss (5 of 5 dropped, block counter 5), while the benign client still had 0\,\% loss. This is the hook through which an IDS or our ML classifier can block a source (\cref{sec:t4-bridge}). When the loader was stopped, the program was detached, as \texttt{ip link show} confirmed.

\screenshot[0.9\linewidth]{demo_phase4_skb.png}{Demo phase~4: flood with the rate limiter in generic mode; 500 TCP segments reach the victim, the benign ping has 0\,\% loss, and the loader shows more than 110,000 drops per second.}
\screenshot[0.9\linewidth]{demo_phase5_maps.png}{Demo phase~5: \texttt{bpftool map show} of \texttt{rl\_state} and the per-CPU counters in \texttt{rl\_stats} (the full dump of \texttt{rl\_state} and \texttt{rl\_config} is in \texttt{logs/demo.log}).}
\screenshot[0.9\linewidth]{demo_phase6_blocklist.png}{Demo phase~6: runtime limit change via \texttt{bpftool map update} and block-listing of the attacker (attacker 100\,\% loss, benign 0\,\% loss).}

\section{Problems and Fixes}\label{sec:t4-problems}
Two bugs in the first demo script cost us time. (1)~\textbf{tcpdump did not terminate.} We first counted the packets reaching the victim with \texttt{tcpdump} in the victim namespace, stopped by \texttt{timeout} with SIGTERM. Already in phase~1 (baseline, no XDP attached) the script hung: tcpdump caught the signal but only acts on it once its blocking read returns, and after the flood had ended no further matching packet arrived, so \texttt{wait} blocked. We replaced it with the victim namespace's own SNMP counter \texttt{TcpInSegs} from \texttt{/proc/net/snmp}, read before and after each round, which also counts exactly what reached the TCP stack. (2)~\textbf{The loader was not stopped.} The loader is started in the background through a shell function, so \verb|$!| is the PID of a subshell. A non-interactive bash starts background jobs with SIGINT ignored, so our SIGINT neither stopped the subshell nor reached the loader, and the program stayed attached. The fix signals the child with SIGTERM (\texttt{pkill -TERM -P \$LOADER\_PID}), which the loader handles by detaching cleanly. A trap removes all namespaces, and with them all XDP attachments, on every exit.

\section{Limitations and Bridge to the ML Model}\label{sec:t4-limits}\label{sec:t4-bridge}
(1)~\textbf{Spoofed sources.} Per-source limiting assumes that the source IP identifies the sender. A SYN flood with random spoofed sources sends one or two packets per address and is never limited; it also churns the LRU map. Real systems add aggregate limits (per destination or port), SYN cookies and signature rules like those of Gatebot/L4Drop \autocite{cf_l4drop}. (2)~\textbf{Window boundary.} A fixed window lets up to $2\times$limit packets pass around a boundary; a token bucket would smooth this. (3)~\textbf{Shared addresses.} Users behind one carrier-grade NAT share one budget; the allow list helps for known gateways. (4)~\textbf{IPv4 only.} IPv6 is passed unfiltered. (5)~\textbf{Ingress only, small test bed.} XDP sees only received packets, and our tests used veth on one host with a single attacker.

The block-list map \texttt{rl\_block} is the interface between \emph{detection} and \emph{enforcement}. When a user-space classifier (our Lab~3 model or a Lab~4 detector) flags a source, it inserts the address with \texttt{bpf\_map\_update\_elem()} or \texttt{bpftool map update name rl\_block key A B C D value 1}; from the next packet on, that source is dropped in the driver, as phase~6 demonstrated. This is the division of labour Cloudflare uses at scale: Gatebot detects in user space, L4Drop enforces in XDP \autocite{cf_gatebot,cf_l4drop}. The classifier's decision latency (\cref{sec:t3-cost}) then matters once per flow, not per packet.
